%% Sect. 2 -- Sects. 2.1-2.6 from augur_AMT_§2 draft v1 (2026-09-19); Sect. 2.7 from the §2.9_3 draft (2026-09-21).
%% 2026-09-29 edit: implementation detail (QR, Eq. 7, Kaufman, relation to DOAS codes, scaling note, unused options)
%% moved to Appendix B5-B6.
\section{Retrieval formulation}
\label{sec:2}

\subsection{Forward model in the extinction domain}
\label{sec:2.1}

The cavity output is converted to an extinction coefficient before any spectral fitting
is attempted. For a sample spectrum $I(\lambda)$ and a reference spectrum
$I_0(\lambda)$ recorded with the cavity filled by zero air,
\begin{equation}
\alpha(\lambda) \;=\;
\Bigl[\underbrace{\tfrac{R_L\,(1-R(\lambda,t))}{d}}_{\textstyle \alpha_{\mathrm{mirror}}}
\;+\; R_L\,\alpha_{\mathrm{Ray}}(\lambda; T_0, P_0)\Bigr]
\frac{I_0(\lambda)-I(\lambda)}{I(\lambda)}
\;-\;\bigl[\alpha_{\mathrm{Ray}}(\lambda; T, P)-\alpha_{\mathrm{Ray}}(\lambda; T_0,P_0)\bigr],
\label{eq:1}
\end{equation}
where $R(\lambda,t)$ is the mirror reflectivity and $d$ the mirror separation. $R_L$ is the ratio of the full mirror
separation to the cavity length actually filled by the sample, excluding the purge-gas volumes at the mirrors, and
$\alpha_{\mathrm{Ray}}$ is the Rayleigh extinction at the temperature and pressure of the sample $(T,P)$ and of the
reference $(T_0,P_0)$. The first bracket converts a relative
intensity change into an absolute extinction; the second removes the Rayleigh difference between the two states. Light-level changes between reference and sample enter $\alpha$ directly, a sensitivity that the iterative
cavity-enhanced DOAS method of \citet{Horbanski_2019} removes.
In $\alpha$, unlike in optical density, the gas amounts enter linearly: the quantity to be fitted is a sum of
reference cross sections with concentration coefficients, a smooth background and the periodic structure the
optics impose.

\subsection{Separable structure}
\label{sec:2.2}

Let $p$ index the $n$ detector pixels of the fitting window and let
\begin{equation*}
x(p) \;=\; \frac{2\,(p-p_{\min})}{p_{\max}-p_{\min}} - 1 \;\in\; [-1,1].
\end{equation*}
The model for the measured extinction is
\begin{equation}
\alpha_{\text{model}}(p)\;=\;
\sum_{i=1}^{m} c_i\,\tau_i(T)\,\sigma_i\bigl(p_i'\bigr)
\;+\;\sum_{j=0}^{q} b_j\,T_j\bigl(x(p)\bigr)
\;+\; a_s \sin(f p) + a_c \cos(f p),
\label{eq:2}
\end{equation}
\begin{equation}
p_i' \;=\; (p-p_c)\,s_i + p_c + \delta_i ,
\label{eq:3}
\end{equation}
with $T_j$ the Chebyshev polynomial of the first kind, $p_c$ a fixed window centre,
and $\tau_i(T)=1+\kappa_i (T-T_{\mathrm{ref}})/100$ the temperature scaling of the
$i$-th cross section. Each reference carries its own wavelength shift $\delta_i$
(pixels) and squeeze $s_i$; in practice one gas is designated the anchor and the
companions are linked to it or fixed, which constrains $\theta$ without changing the structure.
The Chebyshev basis on $x\in[-1,1]$ keeps the background block well conditioned, where a monomial basis would
reach condition numbers that make the background coefficients meaningless and leak into the gas coefficients
through the shared normal equations. The etalon enters as a sine--cosine pair at fixed frequency: since
$A\sin(fp+\varphi)=A\cos\varphi\,\sin(fp)+A\sin\varphi\,\cos(fp)$, the pair $(a_s,a_c)$ spans the same model
space as $(A,\varphi)$ but is linear, so only the frequency $f$ is fixed a priori.

Collecting the $k=m+(q+1)+2$ basis vectors into a design matrix $\mathbf{A}(\theta)$
whose columns depend on the nonlinear vector $\theta=(\delta_i,s_i)$ only through
Eq.~(\ref{eq:3}), and writing $\mathbf{W}=\mathrm{diag}(w_1,\dots,w_n)$ for the pixel weights,
the problem is
\begin{equation}
\min_{\theta,\;\mathbf{c}}\;
\bigl\|\mathbf{W}\bigl(\mathbf{y}-\mathbf{A}(\theta)\,\mathbf{c}\bigr)\bigr\|_2^2,
\qquad \mathbf{y}=\alpha_{\text{meas}} .
\label{eq:4}
\end{equation}
This is a separable nonlinear least-squares problem in the sense of \citet{Golub_1973}:
$k$ linear parameters and, for the operational configuration, two to
six nonlinear ones.

\subsection{Variable projection with an analytic projected Jacobian}
\label{sec:2.3}

For any $\theta$ the inner problem is an ordinary linear least-squares problem, so
$\mathbf{c}$ need never be searched. Substituting the minimiser
$\hat{\mathbf{c}}(\theta)=\mathbf{A}^{+}\mathbf{y}$ leaves the projected residual
\begin{equation}
\mathbf{r}(\theta)\;=\;\bigl(\mathbf{I}-\mathbf{A}(\theta)\mathbf{A}(\theta)^{+}\bigr)\mathbf{y}
\;=\;\mathbf{P}^{\perp}_{\mathbf{A}(\theta)}\,\mathbf{y},
\label{eq:5}
\end{equation}
a function of $\theta$ alone, which we minimise with a bounded
Levenberg--Marquardt/trust-region solver. Its Jacobian is supplied analytically in the full Golub--Pereyra form,
\begin{equation}
\frac{\partial \mathbf{r}}{\partial \theta_k}
\;=\;-\Bigl[\;\underbrace{\mathbf{P}^{\perp}\mathbf{D}_k\,\hat{\mathbf{c}}}_{\text{projection term}}
\;+\;\underbrace{(\mathbf{A}^{+})^{\!\top}\mathbf{D}_k^{\top}\,\mathbf{r}}_{\text{second term}}\;\Bigr],
\qquad
\mathbf{D}_k=\frac{\partial \mathbf{A}}{\partial \theta_k}.
\label{eq:6}
\end{equation}
Retaining the second term makes Eq.~(\ref{eq:6}) the exact derivative rather than the \citet{Kaufman_1975}
approximation, which can therefore be checked against finite differences to round-off (Sect.~\ref{sec:3.2}).
The factorisation, the derivative of the gas columns, the invariance under column scaling and the relation to
existing DOAS implementations, none of which forms this derivative, are given in Appendix~\ref{app:B5}.

\subsection{Time-resolved mirror reflectivity}
\label{sec:2.4}

$R$ in Eq.~(\ref{eq:1}) is determined at intervals from helium and zero-air injections, whose known Rayleigh
extinction fixes the mirror loss \citep{Washenfelder_2008}, and it drifts between them. We therefore carry
$R(\lambda,t)$ as a monotone piecewise-cubic (PCHIP) interpolation \citep{Fritsch_1980} through the injection
knots, evaluated at the timestamp of each spectrum, rather than as one value per campaign or per day. Monotone
interpolation cannot overshoot between knots; the cost, quantified in Sect.~\ref{sec:4.4}, is that it also
cannot represent a turning point that falls between two knots.

\subsection{Uncertainty}
\label{sec:2.5}

With $\hat{\sigma}^2 = \mathbf{r}^{\!\top}\mathbf{r}/(n-P)$, where $P$ counts the
linear columns and the active nonlinear parameters, the conditional covariance of the
linear coefficients is the usual
\begin{equation}
\mathrm{Cov}\bigl(\mathbf{c}\,\big|\,\theta=\hat\theta\bigr)
=\hat{\sigma}^2\bigl(\mathbf{A}_{\mathbf{W}}^{\top}\mathbf{A}_{\mathbf{W}}\bigr)^{-1}.
\label{eq:8}
\end{equation}
Equation~(\ref{eq:8}) is what DOAS retrievals conventionally report. It holds the fitted shift and squeeze at
their optimum: variable projection separates $\theta$ from $\mathbf{c}$ in the solution but not in the
uncertainty, and an uncertain shift passes its uncertainty to the concentrations. We therefore also form the
joint covariance from the full Jacobian of the unprofiled residual,
\begin{equation}
\mathbf{J}=\bigl[\;\mathbf{A}_{\mathbf{W}}\;\big|\;\mathbf{V}\;\bigr],
\qquad
\mathbf{V}_{:,k}=\sum_i \frac{\partial \mathbf{A}_{:,i}}{\partial \theta_k}\,\hat{c}_i,
\qquad
\mathrm{Cov}(\mathbf{c},\theta)=\hat{\sigma}^2\,(\mathbf{J}^{\top}\mathbf{J})^{-1},
\label{eq:9}
\end{equation}
and take the gas block \citep{OLeary_2013}. The joint standard errors are never smaller than the conditional
ones (Appendix~\ref{app:B4}). Both quantities are reported, so that the size of the difference can be read
directly. When $n-P\le 0$ the variance estimate is reported as missing, and the scan is flagged as
underdetermined and retained.

\subsection{Termination state}
\label{sec:2.6}

A residual-and-$\chi^2$ report cannot distinguish a fit that stopped at a parameter bound from one that stopped
because the gradient vanished. We therefore classify and export the termination state of the
nonlinear stage separately from the solver's own success flag: CONVERGED, AT\_BOUND
(a nonlinear parameter rests on its box constraint at exit), STEP\_LIMITED (the
trust-region step was capped), MAX\_NFEV, and FAILED. The classifier is conservative by construction: a
nonlinear parameter within a small tolerance of its box constraint (relative to the box width) is flagged even if
the optimum lies there, and across the solves of one record the least favourable state is kept. It changes no
retrieved value but changes what the record admits (Sects.~\ref{sec:5.2} and \ref{sec:7.3}). Two further implementation matters -- the scaling of $\alpha$ before the solve and the
non-negativity, Tikhonov and robust options, none of which is used operationally -- are described in
Appendix~\ref{app:B6}.

\subsection{What the implementation exposes}
\label{sec:2.9}

Augur makes two questions cheap to ask of a finished retrieval. The first is what happens to the product when a
calibration reference is changed. The structural budget of Sect.~\ref{sec:4} removes one zero-air knot, one
reflectivity knot or one etalon component at a time and re-runs the retrieval, once per knot and then for the pairwise
and triple combinations. It covers the \fieldnum{9034} records of the budget window that have a reflectivity knot within
1\,h (\fieldnum{8847} after bound-terminated fits are excluded). The second is what the objective looks like away from
the converged point: the identifiability map of Sect.~\ref{sec:5} evaluates the projected objective on a grid of 160
points in the wavelength shift for each of twelve spectra, and repeats a coarser version over 165. A fit report
answers neither question; both require re-entering the retrieval with one input altered.

Both questions are loops over data already in memory. The time-resolved references enter the forward model as data,
an array of knots with times interpolated at evaluation, so substituting one is an assignment, and the projected
objective is a callable of the nonlinear parameters that can be evaluated at a chosen shift without converging a fit.
In one process on the analysis workstation, the in-memory budget takes \fieldnum{0.74}\,s per record for both heated channels
and all four perturbations (\fieldnum{74}\,s for a 100-record subset), about \fieldnum{1.9}\,h for the \fieldnum{9034}-record window. A file-based
route that rebuilds the extinction record and refits one channel-day (\fieldnum{1399} records) takes \fieldnum{54}\,s and \fieldnum{90}\,s, about
\fieldnum{4.8}\,min per perturbation for the two channels even if only the affected day is reprocessed. With \fieldnum{157} zero-air and
\fieldnum{54} helium knots in the window, the single-knot removals alone would take about \fieldnum{17}\,h, before the pairwise and
triple combinations. The formulation itself is conventional (Sect.~\ref{sec:3.3}); the contribution of the software is
architectural, in that it brings the quantities of Sects.~\ref{sec:4} and \ref{sec:5} to this cost. The nine-hour time-base error of Sect.~\ref{sec:5.4}, invisible in the residual, could be measured only
by re-running the retrieval with the time base changed.
